% GENERATED FILE. Edit canonical lessons, then run npm run generate:books.
% canonical-source-hash: fnv1a64:3f4f72cd512b563f
% canonical-lessons: ZH-C380-huanjing, ZH-C380-kongqi, ZH-C380-wuran, ZH-C380-diqiu, ZH-C380-huanbao

\chapter{Clean Air}
\label{ch:zh-clean-air}
\hlchaptermodality{380}

\begin{chapteropening}
\textbf{By the end of this chapter:} \emph{I can talk about the environment, the air, pollution, the earth and green ways to travel.}

Complete the last lesson of chapter 380: I can talk about the environment, the air, pollution, the earth and green ways to travel.
\end{chapteropening}

\section[huánjìng]{huánjìng — environment; surroundings}

\label{lesson:ZH-C380-huanjing}



\begin{quote}
Before the new one: say the Mandarin for a card-entry door, then the Mandarin for law and order.
\end{quote}

\subsection*{You'll want to know: huánjìng}
\emph{huánjìng} — \textquotedblleft{}environment; surroundings\textquotedblright{}.

\begin{quote}
\emph{wǒ juéde zhèr de huánjìng bǐ yǐqián hǎo} — I think the environment here is better than before.
\end{quote}

Its characters come later; for now it is a word for the ear.

\begin{cousinweb}
\emph{huánjìng} is built from \emph{huán} 'a ring' and \emph{jìng} 'a border', both inherited from Middle Chinese, \emph{hwaen} and \emph{kjaengX} (Guangyun, in Baxter's transcription), which glosses the first 'a jade ring' and the second 'a boundary': what lies around you up to the edge. The modern sense is not given there.
\end{cousinweb}

\begin{grammarlens}[title={What you've built}]
One more word for giving reasons.
\end{grammarlens}

\subsection*{Your turn}
\begin{itemize}
\raggedright
  \item \emph{Say it:} \emph{huánjìng}
  \item \emph{Say it:} \emph{huánjìng}, once more
  \item \emph{Say it:} \emph{wǒ juéde zhèr de huánjìng bǐ yǐqián hǎo} — I think the environment here is better than before.
  \item \emph{Recall:} say the Mandarin for a card-entry door, then the Mandarin for law and order, then say \emph{huánjìng} again
\end{itemize}

\begin{tcolorbox}[breakable,colback=teal!4,colframe=teal!35!black,title={Before you move on}]
What does huánjìng mean? (\textquotedblleft{}Environment; surroundings\textquotedblright{}.) Say it once more.
\end{tcolorbox}
\section[kōngqì]{kōngqì — air}

\label{lesson:ZH-C380-kongqi}



\begin{quote}
Before the new one: say the Mandarin for law and order, then the Mandarin for the environment.
\end{quote}

\subsection*{You'll want to know: kōngqì}
\emph{kōngqì} — \textquotedblleft{}air\textquotedblright{}.

\begin{quote}
\emph{yīnwèi xiàyǔ le suǒyǐ kōngqì hěn hǎo} — It has rained, so the air is fresh.
\end{quote}

Its characters come later; for now it is a word for the ear.

\begin{cousinweb}
\emph{kōngqì} is built from \emph{kōng} 'empty' and \zh{气} 'breath', both inherited from Middle Chinese, \emph{khuwng} and \emph{khj+jH} (Guangyun, in Baxter's transcription), which glosses the first 'hollow' and the second 'breath'.
\end{cousinweb}

\begin{grammarlens}[title={What you've built}]
One more word for giving reasons.
\end{grammarlens}

\subsection*{Your turn}
\begin{itemize}
\raggedright
  \item \emph{Say it:} \emph{kōngqì}
  \item \emph{Say it:} \emph{kōngqì}, once more
  \item \emph{Say it:} \emph{yīnwèi xiàyǔ le suǒyǐ kōngqì hěn hǎo} — It has rained, so the air is fresh.
  \item \emph{Recall:} say the Mandarin for law and order, then the Mandarin for the environment, then say \emph{kōngqì} again
\end{itemize}

\begin{tcolorbox}[breakable,colback=teal!4,colframe=teal!35!black,title={Before you move on}]
What does kōngqì mean? (\textquotedblleft{}Air\textquotedblright{}.) Say it once more.
\end{tcolorbox}
\section[wūrǎn]{wūrǎn — to pollute}

\label{lesson:ZH-C380-wuran}



\begin{quote}
Before the new one: say the Mandarin for the environment, then the Mandarin for the air we breathe.
\end{quote}

\subsection*{You'll want to know: wūrǎn}
\emph{wūrǎn} — \textquotedblleft{}to pollute\textquotedblright{}.

\begin{quote}
\emph{kāichē huì wūrǎn kōngqì} — Driving pollutes the air.
\end{quote}

Its characters come later; for now it is a word for the ear.

\begin{cousinweb}
\emph{wūrǎn} is built from \emph{wū} 'dirty water' and \emph{rǎn} 'to dye', both inherited from Middle Chinese, \emph{'u} and \emph{nyemX} (Guangyun, in Baxter's transcription), which glosses the first, in another spelling, 'water that stands and does not flow' and the second 'to colour'.
\end{cousinweb}

\begin{grammarlens}[title={What you've built}]
One more word for giving reasons.
\end{grammarlens}

\subsection*{Your turn}
\begin{itemize}
\raggedright
  \item \emph{Say it:} \emph{wūrǎn}
  \item \emph{Say it:} \emph{wūrǎn}, once more
  \item \emph{Say it:} \emph{kāichē huì wūrǎn kōngqì} — Driving pollutes the air.
  \item \emph{Recall:} say the Mandarin for the environment, then the Mandarin for the air we breathe, then say \emph{wūrǎn} again
\end{itemize}

\begin{tcolorbox}[breakable,colback=teal!4,colframe=teal!35!black,title={Before you move on}]
What does wūrǎn mean? (\textquotedblleft{}To pollute\textquotedblright{}.) Say it once more.
\end{tcolorbox}
\section[dìqiú]{dìqiú — the earth; our planet}

\label{lesson:ZH-C380-diqiu}



\begin{quote}
Before the new one: say the Mandarin for the air we breathe, then the Mandarin for to pollute.
\end{quote}

\subsection*{You'll want to know: dìqiú}
\emph{dìqiú} — \textquotedblleft{}the earth; our planet\textquotedblright{}.

\begin{quote}
\emph{dìqiú shì wǒmen dàjiā de jiā} — The earth is home to all of us.
\end{quote}

Its characters come later; for now it is a word for the ear.

\begin{cousinweb}
\emph{dìqiú} is built from \emph{dì} 'the ground' and \emph{qiú} 'a ball', both inherited from Middle Chinese, \emph{dijH} and \emph{gjuw} (Guangyun, in Baxter's transcription), which glosses the first 'the land' and, under another spelling, the second 'a ball of fur that is struck'. The earth as a ball is a later idea.
\end{cousinweb}

\begin{grammarlens}[title={What you've built}]
One more word for giving reasons.
\end{grammarlens}

\subsection*{Your turn}
\begin{itemize}
\raggedright
  \item \emph{Say it:} \emph{dìqiú}
  \item \emph{Say it:} \emph{dìqiú}, once more
  \item \emph{Say it:} \emph{dìqiú shì wǒmen dàjiā de jiā} — The earth is home to all of us.
  \item \emph{Recall:} say the Mandarin for the air we breathe, then the Mandarin for to pollute, then say \emph{dìqiú} again
\end{itemize}

\begin{tcolorbox}[breakable,colback=teal!4,colframe=teal!35!black,title={Before you move on}]
What does dìqiú mean? (\textquotedblleft{}The earth; our planet\textquotedblright{}.) Say it once more.
\end{tcolorbox}
\section[huánbǎo]{huánbǎo — eco-friendly}

\label{lesson:ZH-C380-huanbao}



\begin{quote}
Before the new one: say the Mandarin for to pollute, then the Mandarin for our planet.
\end{quote}

\subsection*{You'll want to know: huánbǎo}
\emph{huánbǎo} — \textquotedblleft{}eco-friendly\textquotedblright{}.

\begin{quote}
\emph{zuò dìtiě bǐ kāichē huánbǎo} — Taking the underground is greener than driving.
\end{quote}

Its characters come later; for now it is a word for the ear.

\begin{cousinweb}
\emph{huánbǎo} is built from \emph{huán}, the first syllable of \emph{huánjìng} 'the environment', and \emph{bǎo} 'to keep', both inherited from Middle Chinese, \emph{hwaen} and \emph{pawX} (Guangyun, in Baxter's transcription), which glosses the second 'to guard': a short form of 'protecting the environment'.
\end{cousinweb}

\begin{grammarlens}[title={What you've built}]
One more word for giving reasons.
\end{grammarlens}

\subsection*{Your turn}
\begin{itemize}
\raggedright
  \item \emph{Say it:} \emph{huánbǎo}
  \item \emph{Say it:} \emph{huánbǎo}, once more
  \item \emph{Say it:} \emph{zuò dìtiě bǐ kāichē huánbǎo} — Taking the underground is greener than driving.
  \item \emph{Recall:} say the Mandarin for to pollute, then the Mandarin for our planet, then say \emph{huánbǎo} again
\end{itemize}

\begin{tcolorbox}[breakable,colback=teal!4,colframe=teal!35!black,title={Before you move on}]
What does huánbǎo mean? (\textquotedblleft{}Eco-friendly\textquotedblright{}.) Say it once more.
\end{tcolorbox}
